\section{AMC 8 1999}
        \begin{enumerate}
        	\item (\textit{AMC 8 1999}) $(6?3) + 4 - (2 - 1) = 5$ To make this statement true, the question mark between the 6 and the 3 should be replaced by


 $\text{(A)} \div \qquad \text{(B)}\ \times \qquad \text{(C)} + \qquad \text{(D)}\ - \qquad \text{(E)}\ \text{None of these}$


 
 \textbf{Answer: }A



\textbf{Solution}

Simplifying the given expression, we get: $(6?3)=2.$

 At this point, it becomes clear that it should be $\boxed{\mathrm{(A)}\ \div}$.

 


	\item (\textit{AMC 8 1999}) What is the degree measure of the smaller angle formed by the hands of a clock at 10 o clock?


 
                    \begin{figure}[H]
                        \centering
                        \includegraphics[width=0.5\textwidth]{Problem/AMC_8_AMC_1999_2_Problem_Figure_1.png}
                    \end{figure}
                    


 $\text{(A)}\ 30 \qquad \text{(B)}\ 45 \qquad \text{(C)}\ 60 \qquad \text{(D)}\ 75 \qquad \text{(E)}\ 90$


 
 \textbf{Answer: }C



\textbf{Solution 1}

At $10:00$, the hour hand will be on the $10$ while the minute hand on the $12$.

 This makes them $\frac{1}{6}$th of a circle apart, and $\frac{1}{6}\cdot360^{\circ}=\boxed{\textbf{(C) } 60}$.

 



\textbf{Solution 2}

We know that the full clock is a circle, and therefore has 360 degrees. Considering that there are $12$ numbers, the distance between one number will be $360\div 12=30$.If the time is $10:00$, then the hour hand will be on $10$, and the minute hand will be on, $12$, making them $2$ numbers apart, so they will be $60$ degrees apart, or $\boxed{\textbf{(C) } 60}$

 


	\item (\textit{AMC 8 1999}) Which triplet of numbers has a sum NOT equal to 1?


 $\text{(A)}\ (1/2,1/3,1/6) \qquad \text{(B)}\ (2,-2,1) \qquad \text{(C)}\ (0.1,0.3,0.6) \qquad \text{(D)}\ (1.1,-2.1,1.0) \qquad \text{(E)}\ (-3/2,-5/2,5)$


 
 \textbf{Answer: }D



\textbf{Solution}

By adding each triplet, we can see that $\boxed{(D)}$ gives us $0$, not $1$, as our sum.

 


	\item (\textit{AMC 8 1999}) The diagram shows the miles traveled by bikers Alberto and Bjorn. After four hours, about how many more miles has Alberto biked than Bjorn?


 
                    \begin{figure}[H]
                        \centering
                        \includegraphics[width=0.5\textwidth]{Problem/AMC_8_AMC_1999_4_Problem_Figure_1.png}
                    \end{figure}
                    


 $\text{(A)}\ 15 \qquad \text{(B)}\ 20 \qquad \text{(C)}\ 25 \qquad \text{(D)}\ 30 \qquad \text{(E)}\ 35$


 
 \textbf{Answer: }A



\textbf{Solution}

After 4 hours, we see that Bjorn biked 45 miles, and Alberto biked 60.  Thus the answer is $60-45=15$ $\boxed{\text{(A)}}$.

 



\textbf{Solution 2}

We see that each dot is $15$ units away from the nearest one above it. So the answer is $\boxed{\text{A}}$.

 


	\item (\textit{AMC 8 1999}) A rectangular garden 60 feet long and 20 feet wide is enclosed by a fence. To make the garden larger, while using the same fence, its shape is changed to a square. By how many square feet does this enlarge the garden?


 $\text{(A)}\ 100 \qquad \text{(B)}\ 200 \qquad \text{(C)}\ 300 \qquad \text{(D)}\ 400 \qquad \text{(E)}\ 500$


 
 \textbf{Answer: }D



\textbf{Solution 1}

We need the same perimeter as a $60$ by $20$ rectangle, but the greatest area we can get. right now the perimeter is $160$. To get the greatest area while keeping a perimeter of $160$, the sides should all be $40$. that means an area of $1600$. Right now, the area is $20 \times 60$ which is $1200$. $1600-1200=400$ which is $\boxed{D}$.

 



\textbf{Solution 2}

We need the same perimeter as a $60$ by $20$ rectangle, but in a square. Right now the perimeter is $160$, since a square has equal sides and $160$ divided by 4 is $40$,  the sides should all be $40$, that means an area of $1600$. Before, the area is $20 \times 60$ which is $1200$. $1600-1200=400$ which is $\boxed{D}$. 

 


	\item (\textit{AMC 8 1999}) Bo, Coe, Flo, Joe, and Moe have different amounts of money. Neither Joe nor Bo has as much money as Flo. Both Bo and Coe have more than Moe. Joe has more than Moe, but less than Bo. Who has the least amount of money?


 $\text{(A)}\ \text{Bo} \qquad \text{(B)}\ \text{Coe} \qquad \text{(C)}\ \text{Flo} \qquad \text{(D)}\ \text{Joe} \qquad \text{(E)}\ \text{Moe}$


 
 \textbf{Answer: }E



\textbf{Solution}

Use logic to solve this problem. You don't actually need to use any equations.

 Neither Jo nor Bo has as much money as Flo. So Flo clearly does not have the least amount of money. Rule out Flo.

 Both Bo and Coe have more than Moe. Rule out Bo and Coe; they clearly do not have the least amount of money.

 Jo has more than Moe. Rule out Jo. 

 The only person who has not been ruled out is Moe. So $\boxed{\text{(E)}}$ is the answer.

 


 


	\item (\textit{AMC 8 1999}) The third exit on a highway is located at milepost 40 and the tenth exit is at milepost 160. There is a service center on the highway located three-fourths of the way from the third exit to the tenth exit. At what milepost would you expect to find this service center?


 $\text{(A)}\ 90 \qquad \text{(B)}\ 100 \qquad \text{(C)}\ 110 \qquad \text{(D)}\ 120 \qquad \text{(E)}\ 130$


 
 \textbf{Answer: }E



\textbf{Solution}

There are $160-40=120$ miles between the third and tenth exits, so the service center is at milepost $40+(3/4)(120) = 40+90=\boxed{\text{(E)}\ 130}$.

 


	\item (\textit{AMC 8 1999}) Six squares are colored, front and back, (R = red, B = blue, O = orange, Y = yellow, G = green, and W = white). They are hinged together as shown, then folded to form a cube. The face opposite the white face is


 
                    \begin{figure}[H]
                        \centering
                        \includegraphics[width=0.5\textwidth]{Problem/AMC_8_AMC_1999_8_Problem_Figure_1.png}
                    \end{figure}
                    


 $\text{(A)}\ \text{B} \qquad \text{(B)}\ \text{G} \qquad \text{(C)}\ \text{O} \qquad \text{(D)}\ \text{R} \qquad \text{(E)}\ \text{Y}$


 
 \textbf{Answer: }A



\textbf{Solution 1}

When G is arranged to be the base, B is the back face and W isthe front face. Thus, $\boxed{\text{(A)}\ B}$ is opposite W.

 



\textbf{Solution 2}

Let Y be the top and fold G, O, and W down.Then $\boxed{\text{(A)}\ B}$ will fold to become the back face and beopposite W.

 



\textbf{Solution 3 (Last Resort)}

If you're out of time or checking your work, I guess you can just rip a net from some paper (no scissors though). Like I said, this is a last resort and not very convenient. $\boxed{\text{(A)}\ B}$ would be the answer.

 


	\item (\textit{AMC 8 1999}) Three flower beds overlap as shown. Bed A has 500 plants, bed B has 450 plants, and bed C has 350 plants. Beds A and B share 50 plants, while beds A and C share 100. The total number of plants is


 
                    \begin{figure}[H]
                        \centering
                        \includegraphics[width=0.5\textwidth]{Problem/AMC_8_AMC_1999_9_Problem_Figure_1.png}
                    \end{figure}
                    


 $\text{(A)}\ 850 \qquad \text{(B)}\ 1000 \qquad \text{(C)}\ 1150 \qquad \text{(D)}\ 1300 \qquad \text{(E)}\ 1450$


 
 \textbf{Answer: }C



\textbf{Solution 1}

Plants shared by two beds have been countedtwice, so the total is $500 + 450 + 350 - 50 - 100 = \boxed{\text{(C)}\ 1150}$.

 



\textbf{Solution 2}

Bed A has $350$ plants it doesn'tshare with B or C. Bed B has $400$ plants it doesn'tshare with A or C. And C has $250$ it doesn't sharewith A or B. The total is $350 + 400 + 250 + 50 + 100 = \boxed{\text{(C)}\ 1150}$ plants.

 


	\item (\textit{AMC 8 1999}) A complete cycle of a traffic light takes 60 seconds. During each cycle the light is green for 25 seconds, yellow for 5 seconds, and red for 30 seconds. At a randomly chosen time, what is the probability that the light will NOT be green?


 $\text{(A)}\ \frac{1}{4} \qquad \text{(B)}\ \frac{1}{3} \qquad \text{(C)}\ \frac{5}{12} \qquad \text{(D)}\ \frac{1}{2} \qquad \text{(E)}\ \frac{7}{12}$


 
 \textbf{Answer: }E



\textbf{Solution 1}


 \[\frac{\text{time not green}}{\text{total time}} = \frac{R + Y}{R + Y + G} = \frac{35}{60} = \boxed{\text{(E)}\ \frac{7}{12}}\]

 



\textbf{Solution 2}

The probability of green is$\frac{25}{60} = \frac{5}{12}$, so the probability of not green is $1-  \frac{5}{12} = \boxed{\text{(E)}\ \frac{7}{12}}$.

 


	\item (\textit{AMC 8 1999}) Each of the five numbers 1, 4, 7, 10, and 13 is placed in one of the five squares so that the sum of the three numbers in the horizontal row equals the sum of the three numbers in the vertical column. The largest possible value for the horizontal or vertical sum is


 
                    \begin{figure}[H]
                        \centering
                        \includegraphics[width=0.5\textwidth]{Problem/AMC_8_AMC_1999_11_Problem_Figure_1.png}
                    \end{figure}
                    


 $\text{(A)}\ 20 \qquad \text{(B)}\ 21 \qquad \text{(C)}\ 22 \qquad \text{(D)}\ 24 \qquad \text{(E)}\ 30$


 
 \textbf{Answer: }D



\textbf{Solution 1}

The largest sum occurs when $13$ is placed in the center. This sum is $13 + 10 + 1 = 13 + 7 + 4 = \boxed{\text{(D)}\ 24}$. Note: Two other common sums, $18$ and $21$, are also possible.

 



\textbf{Solution 2}

Since the horizontal sum equals the vertical sum, twice this sum will be the sumof the five numbers plus the number in the center. When the center number is$13$, the sum is the largest, 
 \[[10 + 4 + 1 + 7 + 2(13)]=2S\\ 48=2S\\ S=\boxed{\text{(D)}\ 24}\]The otherfour numbers are divided into two pairs with equal sums.

 


	\item (\textit{AMC 8 1999}) The ratio of the number of games won to the number of games lost (no ties) by the Middle School Middies is $11/4$. To the nearest whole percent, what percent of its games did the team lose?


 $\text{(A)}\ 24 \qquad \text{(B)}\ 27 \qquad \text{(C)}\ 36 \qquad \text{(D)}\ 45 \qquad \text{(E)}\ 73$


 
 \textbf{Answer: }B



\textbf{Solution}

The ratio means that for every $11$ games won, $4$ are lost, so the team has won $11x$ games, lost $4x$ games, and played $15x$ games for some positive integer $x$. The percentage of games lost is just $\dfrac{4x}{15x}\times100=\dfrac{4}{15}\times 100=26.\overline{6}\%\approx\boxed{\text{(B)}\ 27\%}$

 \textbf{Alternate Solution} The Won/Lost ratio is 11/4 so, for some number $N$, the team won $11N$ games and lost $4N$ games. Thus, the team played $15N$ games and the fraction of games lost is $\dfrac{4N}{15N}=\dfrac{4}{15}\approx 0.27=\boxed{27\%}.$-Clara Garza

 


	\item (\textit{AMC 8 1999}) The average age of the 40 members of a computer science camp is 17 years. There are 20 girls, 15 boys, and 5 adults. If the average age of the girls is 15 and the average age of the boys is 16, what is the average age of the adults?


 $\text{(A)}\ 26 \qquad \text{(B)}\ 27 \qquad \text{(C)}\ 28 \qquad \text{(D)}\ 29 \qquad \text{(E)}\ 30$


 
 \textbf{Answer: }C



\textbf{Solution}

First, find the total amount of the girl's ages and add it to the total amount of the boy's ages. It equals $(20)(15)+(15)(16)=540$. The total amount of everyone's ages can be found from the average age, $17\cdot40=680$. Then you do $680-540=140$ to find the sum of the adult's ages. The average age of an adult is divided among the five of them, $140\div5=\boxed{\text{(C)}\ 28}$.

 


	\item (\textit{AMC 8 1999}) In trapezoid $ABCD$, the sides $AB$ and $CD$ are equal. The perimeter of $ABCD$ is


 
                    \begin{figure}[H]
                        \centering
                        \includegraphics[width=0.5\textwidth]{Problem/AMC_8_AMC_1999_14_Problem_Figure_1.png}
                    \end{figure}
                    


 $\text{(A)}\ 27 \qquad \text{(B)}\ 30 \qquad \text{(C)}\ 32 \qquad \text{(D)}\ 34 \qquad \text{(E)}\ 48$


 
 \textbf{Answer: }D



\textbf{Solution}


                    \begin{figure}[H]
                        \centering
                        \includegraphics[width=0.5\textwidth]{Solution/AMC_8_AMC_1999_14_Solution_Figure_1.png}
                    \end{figure}
                    

 There is a rectangle present, with both horizontal bases being $8$ units in length. The excess units on the bottom base must then be $16-8=8$. The fact that $AB$ and $CD$ are equal in length indicate, by the Pythagorean Theorem, that these excess lengths are equal. There are two with a total length of $8$ units, so each is $4$ units. The triangle has a hypotenuse of $5$, because the triangles are $3-4-5$ right triangles. So, the sides of the trapezoid are $8$, $5$, $16$, and $5$. Adding those up gives us the perimeter, $8 + 5 + 16 + 5 = 13 + 21 = \boxed{\text{(D)}\ 34}$ units.

 


	\item (\textit{AMC 8 1999}) Bicycle license plates in Flatville each contain three letters. The first is chosen from the set {C,H,L,P,R}, the second from {A,I,O}, and the third from {D,M,N,T}. 


 When Flatville needed more license plates, they added two new letters. The new letters may both be added to one set or one letter may be added to one set and one to another set. What is the largest possible number of ADDITIONAL license plates that can be made by adding two letters?


 $\text{(A)}\ 24 \qquad \text{(B)}\ 30 \qquad \text{(C)}\ 36 \qquad \text{(D)}\ 40 \qquad \text{(E)}\ 60$


 
 \textbf{Answer: }D



\textbf{Solution 1 (Complementary Counting)}

There are currently $5$ choices for the first letter, $3$ choices for the second letter, and $4$ choices for the third letter, for a total of $5 \cdot 3 \cdot 4 = 60$ license plates.

 Adding $2$ letters to the start gives $7\cdot 3 \cdot 4 = 84$ plates.

 Adding $2$ letters to the middle gives $5 \cdot 5 \cdot 4 = 100$ plates.

 Adding $2$ letters to the end gives $5 \cdot 3 \cdot 6 = 90$ plates.

 Adding a letter to the start and middle gives $6 \cdot 4 \cdot 4 = 96$ plates.

 Adding a letter to the start and end gives $6 \cdot 3 \cdot 5 = 90$ plates.

 Adding a letter to the middle and end gives $5 \cdot 4 \cdot 5 = 100$ plates.

 You can get at most $100$ license plates total, giving an additional $100 - 60 = 40$ plates, making the answer $\boxed {D}$

 



\textbf{Solution 2}

Using the same logic as above, the number of combinations of plates is simply the product of the size of each set of letters.

 In general, when three numbers have the same fixed sum, their product will be maximal when they are as close together as possible.  This is a 3D analogue of the fact that a rectangle with fixed perimeter maximizes its area when the sides are equal (ie when it becomes a square).  In this case, no matter where you add the letters, there will be $5 + 3 + 4 + 2 = 14$ letters in total.  If you divide them as evenly as possible among the three groups, you get $5, 5, 4$, which is a possible situation.

 As before, the answer is $5 \cdot 5 \cdot 4 - 5 \cdot 3 \cdot 4 = 40$, and the correct choice is $\boxed{D}$

 



\textbf{Solution 3}

Before new letters were added, five different letters could have been chosen for the first position, three for the second, and four for the third. This means that $5\cdot 3\cdot 4=60$ plates could have been made.

 If two letters are added to the second set, then $5\cdot 5\cdot 4=100$ plates can be made. If one letter is added to each of the second and third sets, then $5\cdot 4\cdot 5=100$ plates can be made. None of the other four ways to place the two letters will create as many plates. So, $100-60=\boxed{40}$ ADDITIONAL plates can be made.So the correct choice is $\boxed{D}$

 


	\item (\textit{AMC 8 1999}) Tori's mathematics test had 75 problems: 10 arithmetic, 30 algebra, and 35 geometry problems. Although she answered 70\% of the arithmetic, 40\% of the algebra, and 60\% of the geometry problems correctly, she did not pass the test because she got less than 60\% of the problems right. How many more problems would she have needed to answer correctly to earn a 60\% passing grade?


 $\text{(A)}\ 1 \qquad \text{(B)}\ 5 \qquad \text{(C)}\ 7 \qquad \text{(D)}\ 9 \qquad \text{(E)}\ 11$


 
 \textbf{Answer: }B



\textbf{Solution}

First, calculate how many of each type of problem she got right:

 Arithmetic:  $70\% \cdot 10 = 0.70 \cdot 10 = 7$

 Algebra:  $40\% \cdot 30 = 0.40 \cdot 30 = 12$

 Geometry:  $60\% \cdot 35 = 0.60 \cdot 35 =  21$

 Altogether, Tori answered $7 + 12 + 21 = 40$ questions correct.To get a $60\%$ on her test overall, she needed to get $60\% \cdot 75 = 0.60 \cdot 75 = 45$ questions right.

 Therefore, she needed to answer $45 - 40 = 5$ more questions to pass, so the correct answer is $\boxed{(B) 5}$

 


	\item (\textit{AMC 8 1999}) Problems 17, 18, and 19 refer to the following:


 At Central Middle School the 108 students who take the 1999 AMC 8 meet in the evening to talk about problems and eat an average of two cookies apiece. Walter and Gretel are baking Bonnie's Best Bar Cookies this year. Their recipe, which makes a pan of 15 cookies, lists these items: $1\frac{1}{2}$ cups flour, $2$ eggs, $3$ tablespoons butter, $\frac{3}{4}$ cups sugar, and $1$ package of chocolate drops. They will make only full recipes, not partial recipes.


 Walter can buy eggs by the half-dozen. How many half-dozens should he buy to make enough cookies? (Some eggs and some cookies may be left over.)


 $\text{(A)}\ 1 \qquad \text{(B)}\ 2 \qquad \text{(C)}\ 5 \qquad \text{(D)}\ 7 \qquad \text{(E)}\ 15$


 
 \textbf{Answer: }C



\textbf{Solution}

If $108$ students eat $2$ cookies on average, there will need to be $108\cdot 2 = 216$ cookies.  There are $15$ cookies per pan, meaning there needs to be $\frac{216}{15} = 14.4$ pans.  However, since half-recipes are forbidden, we need to round up and make $\lceil \frac{216}{15}\rceil = 15$ pans.

 $1$ pan requires $2$ eggs, so $15$ pans require $2\cdot 15 = 30$ eggs.   Since there are $6$ eggs in a half dozen, we need $\frac{30}{6} = 5$ half-dozens of eggs, and the answer is $\boxed{C}$

 


	\item (\textit{AMC 8 1999}) Problems 17, 18, and 19 refer to the following:


 At Central Middle School the 108 students who take the AMC 8 meet in the evening to talk about problems and eat an average of two cookies apiece. Walter and Gretel are baking Bonnie's Best Bar Cookies this year. Their recipe, which makes a pan of 15 cookies, lists these items: $1\frac{1}{2}$ cups flour, $2$ eggs, $3$ tablespoons butter, $\frac{3}{4}$ cups sugar, and $1$ package of chocolate drops. They will make only full recipes, not partial recipes.


 They learn that a big concert is scheduled for the same night and attendance will be down 25\%. How many recipes of cookies should they make for their smaller party?


 $\text{(A)}\ 6 \qquad \text{(B)}\ 8 \qquad \text{(C)}\ 9 \qquad \text{(D)}\ 10 \qquad \text{(E)}\ 11$


 
 \textbf{Answer: }E



\textbf{Solution 1}

If $108$ students eat $2$ cookies on average, there will need to be $108\cdot 2 = 216$ cookies.  But with the smaller attendance, you will only need $100\% - 25\% = 75\%$ of these cookies, or $75\% \cdot 216 = 0.75\cdot 216 = 162$ cookies.

 $162$ cookies requires $\frac{162}{15} = 10.8$ batches.  However, since half-batches are forbidden, we must round up to get $\left\lceil \frac{162}{15} \right\rceil = 11$ batches, and the correct answer is $\boxed{E}$.

 



\textbf{Solution 2}

If there were $108$ students before, with the $25\%$ of students missing, there will be $75\%$ of $108$ students left.  This is $75\% \cdot 108 =  0.75 \cdot 108 = 81$ students.  These students eat $81 \cdot 2 = 162$ cookies.  Follow the logic of the second paragraph above to find that there needs to be $11$ batches, and the correct answer is $\boxed{E}$.

 


 


	\item (\textit{AMC 8 1999}) Problems 17, 18, and 19 refer to the following:


 At Central Middle School the 108 students who take the AMC 8 meet in the evening to talk about problems and eat an average of two cookies apiece. Walter and Gretel are baking Bonnie's Best Bar Cookies this year. Their recipe, which makes a pan of 15 cookies, lists these items: $1\frac{1}{2}$ cups flour, $2$ eggs, $3$ tablespoons butter, $\frac{3}{4}$ cups sugar, and $1$ package of chocolate drops. They will make only full recipes, not partial recipes.


 The drummer gets sick. The concert is cancelled. Walter and Gretel must make enough pans of cookies to supply 216 cookies. There are 8 tablespoons in a stick of butter. How many sticks of butter will be needed? (Some butter may be left over, of course.)


 $\text{(A)}\ 5 \qquad \text{(B)}\ 6 \qquad \text{(C)}\ 7 \qquad \text{(D)}\ 8 \qquad \text{(E)}\ 9$


 
 \textbf{Answer: }B



\textbf{Solution}

For $216$ cookies, you need to make $\frac{216}{15} = 14.4$ pans.  Since fractional pans are forbidden, round up to make $\lceil \frac{216}{15} \rceil = 15$ pans.

 There are $3$ tablespoons of butter per pan, meaning $3 \cdot 15 = 45$ tablespoons of butter are required for $15$ pans.

 Each stick of butter has $8$ tablespoons, so we need $\frac{45}{8} = 5.625$ sticks of butter.  However, we must round up again because partial sticks of butter are forbidden!  Thus, we need $\lceil \frac{45}{8} \rceil = 6$ sticks of butter, and the answer is $\boxed{B}$.

 


	\item (\textit{AMC 8 1999}) Figure 1 is called a "stack map." The numbers tell how many cubes are stacked in each position. Fig. 2 shows these cubes, and Fig. 3 shows the view of the stacked cubes as seen from the front.


 Which of the following is the front view for the stack map in Fig. 4?


 
                    \begin{figure}[H]
                        \centering
                        \includegraphics[width=0.5\textwidth]{Problem/AMC_8_AMC_1999_20_Problem_Figure_1.png}
                    \end{figure}
                    


 
 



 
                    \begin{figure}[H]
                        \centering
                        \includegraphics[width=0.5\textwidth]{Problem/AMC_8_AMC_1999_20_Problem_Figure_2.png}
                    \end{figure}
                    


 
 \textbf{Answer: }B



\textbf{Solution}

The third view is a direct, head-on view of the cubes.  Thus, you will only see the highest (or, in these cases, higher) tower in each up-down column.  For figure $4$:

 The highest tower in the first up-down column is $2$ in the upper-left box.

 The highest tower in the second up-down column is the $3$ in the lower-middle box.

 The highest tower in the third up-down column is the $4$ in the upper-right box.

 Thus, the head-on view of this tower should have $2$ boxes on the left, $3$ in the middle, and $4$ on the right.

 Diagram $\boxed{B}$ shows this description.

 As trivia, option C shows the stack from the point of view of an observer on the right, facing towards the left.

 


 


 


	\item (\textit{AMC 8 1999}) The degree measure of angle $A$ is


 
                    \begin{figure}[H]
                        \centering
                        \includegraphics[width=0.5\textwidth]{Problem/AMC_8_AMC_1999_21_Problem_Figure_1.png}
                    \end{figure}
                    


 $\text{(A)}\ 20 \qquad \text{(B)}\ 30 \qquad \text{(C)}\ 35 \qquad \text{(D)}\ 40 \qquad \text{(E)}\ 45$


 
 \textbf{Answer: }B



\textbf{Solution 1}

Angle-chasing using the small triangles:

 Use the line below and to the left of the $110^\circ$ angle to find that the rightmost angle in the small lower-left triangle is $180 - 110 = 70^\circ$.

 Then use the small lower-left triangle to find that the remaining angle in that triangle is $180 - 70 - 40 = 70^\circ$.

 Use congruent vertical angles to find that the lower angle in the smallest triangle containing $A$ is also $70^\circ$.  

 Next, use line segment $AB$ to find that the other angle in the smallest triangle containing $A$ is  $180 - 100 = 80^\circ$.

 The small triangle containing $A$ has a $70^\circ$ angle and an $80^\circ$ angle.  The remaining angle must be $180 - 70 - 80 = \boxed{30^\circ, B}$

 



\textbf{Solution 2 (Quick and Easy, using larger triangles)}

The third angle of the triangle containing the $100^\circ$ angle and the $40^\circ$ angle is $180^\circ - 100^\circ - 40^\circ = 40^\circ$. It follows that $A$ is the third angle of the triangle consisting of the found $40^\circ$ angle and the given $110^\circ$ angle. Thus, $A$ is a $180^\circ - 110^\circ - 40^\circ = 30^\circ$ angle, and so the answer is $\boxed{30^\circ, \textbf{B}}$.

 


	\item (\textit{AMC 8 1999}) In a far-off land three fish can be traded for two loaves of bread and a loaf of bread can be traded for four bags of rice. How many bags of rice is one fish worth?


 $\text{(A)}\ \frac{3}{8} \qquad \text{(B)}\ \frac{1}{2} \qquad \text{(C)}\ \frac{3}{4} \qquad \text{(D)}\ 2\frac{2}{3} \qquad \text{(E)}\ 3\frac{1}{3}$


 
 \textbf{Answer: }D



\textbf{Solution}

Let $f$ represent one fish, $l$ a loaf of bread, and $r$ a bag of rice. Then:$3f=2l$, $l=4r$

 Substituting $l$ from the second equation back into the first gives us $3f=8r$. So each fish is worth $\frac{8}{3}$ bags of rice, or $2 \frac{2}{3}\Rightarrow \boxed{D}$.

 


	\item (\textit{AMC 8 1999}) Square $ABCD$ has sides of length 3. Segments $CM$ and $CN$ divide the square's area into three equal parts. How long is segment $CM$?


 
                    \begin{figure}[H]
                        \centering
                        \includegraphics[width=0.5\textwidth]{Problem/AMC_8_AMC_1999_23_Problem_Figure_1.png}
                    \end{figure}
                    


 $\text{(A)}\ \sqrt{10} \qquad \text{(B)}\ \sqrt{12} \qquad \text{(C)}\ \sqrt{13} \qquad \text{(D)}\ \sqrt{14} \qquad \text{(E)}\ \sqrt{15}$


 
 \textbf{Answer: }C



\textbf{Solution}

Since the square has side length $3$, the area of the entire square is $9$.

 The segments divide the square into 3 equal parts, so the area of each part is $9 \div 3 = 3$.

 Since $\triangle CBM$ has area $3$ and base $CB = 3$, using the area formula for a triangle:

 $A_{tri} = \frac{1}{2}bh$

 $3 = \frac{1}{2}3h$

 $h = 2$

 Thus, height $BM = 2$.

 Since $\triangle CBM$ is a right triangle, $CM = \sqrt{BM^2 + BC^2} = \sqrt{2^2 + 3^2} = \boxed{\text{(C)}\ \sqrt{13}}$.

 


 



\textbf{Solution 2}

Connect $AC$, $S_\triangle AMC=S_\triangle ANC$. To satisfied the three area is equal, we have $2S_\triangle AMC=S_\triangle BMC$, $2S_\triangle ANC=S_\triangle DNC$. Thus, $AM=AN=\frac{1}{2}BM=\frac{1}{2}ND=1$. $BM=2,BC=3,MC=\boxed{\text{(C)}\ \sqrt{13}}$.

 


	\item (\textit{AMC 8 1999}) When $1999^{2000}$ is divided by $5$, the remainder is 


 $\text{(A)}\ 0 \qquad \text{(B)}\ 1 \qquad \text{(C)}\ 2 \qquad \text{(D)}\ 3 \qquad \text{(E)}\ 4$


 
 \textbf{Answer: }B



\textbf{Solution 1}

Note that the units digits of the powers of 9 have a pattern: $9^1 = {\bf 9}$, $9^2 = 8{\bf 1}$, $9^3 = 72{\bf 9}$, $9^4 = 656{\bf 1}$, and so on. Since all natural numbers with the same last digit have the same remainder when divided by 5, the entire number doesn't matter, just the last digit. For even powers of $9$, the number ends in a $1$. Since the exponent is even, the final digit is $1$. Note that all natural numbers that end in $1$ have a remainder of $1$ when divided by $5$. So, our answer is $\boxed{\text{(B)}\ 1}$.

 



\textbf{Solution 2}

Write $1999$ as $2000-1$. We are taking $(2000-1)^{2000} \mod{10}$. Using the binomial theorem, we see that ALL terms in this expansion are divisible by $2000$, which is in turn divisible by $5$, except for the very last term, which is just $(-1)^{2000}$. This is clear because the binomial expansion is just choosing how many $2000$s and how many $-1$s there are for each term. Using this, we can take the entire polynomial $\mod{10}$, which leaves just $(-1)^{2000}=\boxed{\text{(B)}\ 1}$.

 



\textbf{Solution 3}

As $1999 \equiv -1 \pmod{5}$, we have $1999^{2000} \equiv (-1)^{2000} \equiv 1 \pmod{5}$. Thus, the answer is $\boxed{\text{(B)}\ 1}$.

 



\textbf{Solution 4}

A sum/product of any two natural numbers has the same remainder, when divided by $5$, as the sum/product of their remainders. Thus, we can use the basic definition of an exponent and view the problem as $1999 \cdot 1999 \cdot \cdot \cdot 1999$. 

 Using the fact stated in the first sentence, we see that the remainder of $1999$, when divided by $5$, is $4$. The problem can now be written viewed as finding the remainder of $4^{2000}$ when it is divided by $5$, which is already much simpler. 

 Now, toying with the simplified problem a little, see notice that powers of $4$ alternate from ending in $4$ and $6$. We notice that even powers of $4$ always end in $6$, and also that $2000$ is even. Thus $4^{2000}$ must end in $6$, which, when divided by $5$ gives a remainder of $\boxed{\text{(B)}\ 1}$.

 



\textbf{Solution 5}

By Fermat's Little Theorem, $1999^{2000} \equiv 1999^{2000\mod (5-1)} \equiv 1999^0 \equiv 1 \pmod 5$. Thus the answer is $\boxed{\text{(B)}\ 1}$.

 


	\item (\textit{AMC 8 1999}) Points $B$, $D$, and $J$ are midpoints of the sides of right triangle $ACG$. Points $K$, $E$, $I$ are midpoints of the sides of triangle $JDG$, etc. If the dividing and shading process is done 100 times (the first three are shown) and $AC=CG=6$, then the total area of the shaded triangles is nearest


 
                    \begin{figure}[H]
                        \centering
                        \includegraphics[width=0.5\textwidth]{Problem/AMC_8_AMC_1999_25_Problem_Figure_1.png}
                    \end{figure}
                    


 $\text{(A)}\ 6 \qquad \text{(B)}\ 7 \qquad \text{(C)}\ 8 \qquad \text{(D)}\ 9 \qquad \text{(E)}\ 10$


 
 \textbf{Answer: }A



\textbf{Solution 1}

Since $\triangle FGH$ is fairly small relative to the rest of the diagram, we can make an underestimate by using the current diagram.  All triangles are right-isosceles triangles.

 
 \[CD = \frac {CG}{2} = 3, DE = \frac{CD}{2} = \frac{3}{2}, EF = \frac{DE}{2} = \frac{3}{4}\]
 \[CB = CD = 3, DK = DE = \frac{3}{2}, EL = EF = \frac{3}{4}\]
 \[[CBD] = \frac{1}{2}3^2 = \frac{9}{2}\]
 \[[DKE] = \frac{1}{2}(\frac{3}{2})^2 = \frac{9}{8}\]
 \[[ELF] = \frac{1}{2}(\frac{3}{4})^2 = \frac{9}{32}\]

 The sum of the shaded regions is $\frac{9}{2} + \frac{9}{8} + \frac{9}{32} = \frac{189}{32} \approx 5.9$

 $5.9$ is an underestimate, as some portion (but not all) of $\triangle FGH$ will be shaded in future iterations.  

 If you shade all of $\triangle FGH$, this will add an additional $\frac{9}{32}$ to the area, giving $\frac{198}{32} \approx 6.2$, which is an overestimate.

 Thus, $\boxed{\text{(A)}\ 6}$ is the only answer that is both over the underestimate and under the overestimate.

 



\textbf{Solution 2}

In iteration $1$, congruent triangles $\triangle ABJ,  \triangle BDJ,$ and $\triangle BDC$ are created, with one of them being shaded.

 In iteration $2$, three more congruent triangles are created, with one of them being shaded.

 As the process continues indefinitely, in each row, $\frac{1}{3}$ of each triplet of new congruent triangles will be shaded.  The "fourth triangle" at the top ($\triangle FGH$ in the diagram) will gradually shrink, 

 leaving about $\frac{1}{3}$ of the area shaded.  This means $\frac{1}{3}\left(\frac{1}{2}6\cdot 6\right) = 6$ square units will be shaded when the process goes on indefinitely, giving $\boxed{A}$.

 



\textbf{Solution 3}

Using Solution 1 as a template, note that the sum of the areas forms a \href{https://artofproblemsolving.com/wiki/index.php?title=Geometric\_series}{geometric series}:

 $\frac{9}{2} + \frac{9}{8} + \frac{9}{32} + \frac{9}{128} + ...$

 This is the sum of a geometric series with first term $a_1 = \frac{9}{2}$ and common ratio $r = \frac{1}{4}$ This is the easiest way to do this problem.

 The sum of an infinite geometric series with $|r|<1$ is shown by the formula. $S_{\infty} = \frac{a_1}{1 - r}$ Insert the values to get $\frac{\frac{9}{2}}{1 - \frac{1}{4}} = \frac{9}{2}\cdot\frac{4}{3} = 6$, giving an answer of $\boxed{A}$.

 



\textbf{Solution 4}

Find the area of the bottom triangle, which is 4.5. Notice that the area of the triangles is divided by 4 every time. 4.5*5/4\approx5.7, and 5.7+(1/4)\approx5.9. We can clearly see that the sum is approaching answer choice $\boxed{A}$.

 



\textbf{Solution 5}

The area of $DBC$ is clearly $\frac{1}{3}$ of $AJDC$, the area of $DEK$ is $\frac{1}{3}$ of $JIED$, if the progress is going to infinity, the shaded triangles will be $\frac{1}{3}$ of the triangle $ACG$. However, 100 times is much enough. The answer is $\frac{1}{3}\times 6\times 6\times \frac{1}{2} = \boxed{(A)6}$.

 



\textbf{Solution 6}

Because all of the triangles are isosceles, this gives us the area of one triangle with side length $s$, $A = \frac{1}{2}s^2$. Logically, $s$ keeps on getting $\frac{1}{2}$ as small as the previous one ($3, \frac{3}{2}, \frac{3}{4}$, etc) Meaning for triangle number $n, s = 3(\frac{1}{2})^n$. Putting this into our original equation, $A = \frac{1}{2}(3(\frac{1}{2})^n)^2 = \frac{9}{2}(\frac{1}{2})^{2n} = 9(\frac{1}{2})^{2n+1}$ This means the total area is $9\sum^{\infty}_{n=1} \frac{1}{2}^{2n+1}$, using exponent laws, we can simplify this to $\frac{9}{2} \sum^{\infty}_{n=1} \frac{1}{4}^n$.

 The formula for a geometric series like this is $\sum^{\infty}_{n=1} a*r^n = \frac{a}{1-r}$. Using this, we get $\frac{\frac{9}{2}}{\frac{4}{4}-\frac{1}{4}} = \frac{9}{2}\times\frac{4}{3} = 3\times2 = \boxed{(A)6}$.

 



\textbf{Solution 7}

Notice how you can fit $\triangle KDE$ into the top right corner of square $BCDJ$? Repeat this process by putting smaller triangles into the corner of the square and eventually, you should see that this converges to just $\frac{1}{4}$ of the square. $\triangle BCD$ is $\frac{1}{2}$ of the square and rest of the triangles make $\frac{1}{4}$ which totals to $\frac{3}{4}$ of the square. Square $BCDJ$ has side length $3$ so it's area is $9$, and $\frac{3}{4}$ of $9$ is $\boxed{(A)6}$.

 



\textbf{Solution 8}

As in Solution 3, the area form the geometric series $\frac{9}{2} + \frac{9}{8} + \frac{9}{32} + \frac{9}{128} + ...$. If this is $x$, than 4$x$ is 18 + $\frac{9}{2} + \frac{9}{8} + \frac{9}{32} + \frac{9}{128} + ...$.Than we take away $x$ to get 3$x$ = 18. When we solve for $x$ we get $x$ = 6, so the correct answer is $\boxed{A}$.

 


 


\end{enumerate}